\chapter{Caste and Gender}

\section{Rudolph \& Rudolph: Three Forms of Mobilization}

\begin{KeyConcept}
\begin{itemize}
\item \textbf{Lloyd and Susanne Rudolph} hold that the relation between caste and politics in India can be explained through \textbf{three forms of mobilization}:
\item 1. \textbf{Vertical mobilization} --- traditional \textbf{upper-caste elites mobilise the (largely low-caste) masses} (e.g. the national movement: Gandhi, Patel, the Indian National Congress). Vertical mobilization means \textbf{the status quo is maintained} --- nothing really changes.
\item 2. \textbf{Horizontal mobilization} --- \textbf{similar caste groups are mobilised by individuals of the same caste} (what M. N. Srinivas called \textbf{horizontal solidarity}) --- e.g. the \textbf{Mahad Satyagraha}, Ambedkar's \textbf{temple-entry movement}, the \textbf{Dalit Panther} movement, the \textbf{Kang/Kaung Satyagraha}.
\item 3. \textbf{Differential mobilization} --- the masses are mobilised by a \textbf{political party} (not a caste association), cutting across castes --- e.g. the \textbf{Bahujan Samajwadi Party} (it has Brahmin, OBC, Shudra and ST candidates, even though the name is 'Bahujan'). The BSP's roots go back to \textbf{BAMCEF} and \textbf{DS4 (Dalit Shoshit Samaj Sangharsh Samiti)} --- caste associations that grew into a political party.
\end{itemize}
\end{KeyConcept}

\begin{KeyConcept}
\begin{itemize}
\item Caste associations (e.g. the All-India Yadav Welfare Association) are \textbf{horizontal} mobilization --- same-caste people organising same-caste groups.
\item \textbf{More and more differential mobilization = more democratization} (people affiliate to \textbf{ideology/development}, not caste); more vertical mobilization = nothing changes.
\item The Rudolphs' conclusion: \textbf{prima facie it seems caste associations and their growth counter the modernizing process unleashed by colonialism; however, in reality, caste associations turned out to be the agents of modernity in India. Earlier we had casteization of politics; now we have politicization of caste.}
\end{itemize}
\end{KeyConcept}

\section{Anil Bhatt: Four Ways Caste Relates to Politics}

\begin{KeyConcept}
\begin{itemize}
\item \textbf{Anil Bhatt} (drawing on extensive surveys, e.g. by \textbf{CSDS} --- the 'psephologists'/political analysts like \textbf{Yogendra Yadav} and \textbf{Prashant Kishor}) analysed the relation between caste and politics in \textbf{four ways}:
\item 1. \textbf{Political interest} --- \textbf{high-status castes have higher levels of political interest} (they have always been at the top and want to maintain that hegemony in secular status too).
\item 2. \textbf{Political awareness} --- high castes also have higher awareness.
\item 3. \textbf{Political identification (partisanship)} --- \textbf{no relationship was found between caste status and partisanship}: persons of low caste are almost as likely to support a party of the high-caste status group (you support a party on its agenda/development, not on your caste).
\item 4. \textbf{Political influence} --- \textbf{high-status castes are not necessarily the influential castes}: middle and low castes dominate in many villages (the \textbf{dominant caste} --- Patidars in Gujarat, Yadavs in UP, Reddys/Kammas in Andhra).
\end{itemize}
\end{KeyConcept}

\begin{KeyConcept}
\begin{itemize}
\item \textbf{Before independence}, interest, awareness, identification and influence \textbf{intersected with each other}; \textbf{after independence they no longer do} --- so the undifferentiated India became more differentiated. One caste may have high interest but low influence, and vice versa.
\item This means \textbf{caste is getting democratized} (not the reverse).
\end{itemize}
\end{KeyConcept}

\section{Hierarchy vs Difference: Which Scholars Said What}

\begin{KeyConcept}
\begin{itemize}
\item On the question '\textbf{which is more significant --- the principle of hierarchy or the principle of difference in inter-caste relations?}' the scholars divide:
\item \textbf{Principle of hierarchy} --- the Indological fraternity: \textbf{Ghurye} (a \textbf{definite, fixed, rigid hierarchy} --- no change) and \textbf{Dumont} (a \textbf{rigid hierarchy} of purity/pollution --- nothing nebulous; the pure stay pure, the impure stay impure).
\item \textbf{Principle of difference} --- \textbf{Srinivas} (a \textbf{nebulous hierarchy} --- more difference than hierarchy: each caste is different, with claims and counter-claims to superiority) and \textbf{Beteille} (no fixed hierarchy --- Brahmins are not homogeneous but segmented into status groups, so hierarchy depends on context).
\item The line to recall: \textbf{'in the social world, there are certain social forms which are only different, but differences tend to get hierarchized.'} The discussion concludes with \textbf{Dipankar Gupta}.
\end{itemize}
\end{KeyConcept}

\section{Caste Requires Endogamy: Surpanakha and the Smritis}

\begin{ExampleBox}
\begin{itemize}
\item The \textbf{Surpanakha} example: when Surpanakha (a character in the Ramayana) proposed to Rama, he refused, saying '\textbf{I am Kshatriya and you are Brahmin --- we cannot marry}.'
\item Sociological analysis: \textbf{caste requires endogamy for its own growth} --- caste is a \textbf{social fact}, external to the individual and coercive even over Rama. Rama is the \textbf{spokesperson of the caste structure}: it is the caste which regulates the individual.
\end{itemize}
\end{ExampleBox}

\begin{KeyConcept}
\begin{itemize}
\item What restricts our choices is a \textbf{written text} --- the \textbf{Dharma Shastra} (just as \textbf{Sharia} regulates the everyday life of Muslims). The \textbf{Vedas} (Rig, Sama, Yajur, Atharva) are forms of \textbf{knowledge}, not regulators of daily life.
\item \textbf{Shruti} = 'heard' (revealed); \textbf{Smriti} = 'remembered' (human-composed). The \textbf{Yajnavalkya Smriti} says one should marry a woman \textbf{whose virginity is intact, not previously wed by another, dear to one's heart, of the same varna, not a sapinda, younger than oneself, not diseased, and not from the same pravara or gotra}.
\item The \textbf{upper-caste male view} is reproduced through the Yajnavalkya Smriti and legitimised by the Dharma Shastra: you are simply an \textbf{agent through which caste reproduces itself}.
\end{itemize}
\end{KeyConcept}

\section{Uma Chakravarti: Marriage, Sexuality and the Caste System}

\begin{KeyConcept}
\begin{itemize}
\item \textbf{The relation between caste and gender is historical and dynamic.}
\item \textbf{Uma Chakravarti}: \textbf{the structure of marriage and sexuality is the foundation of the caste system.} Marriage can be \textbf{endogamous, hypergamous or hypogamous} (all judged from the woman's perspective) --- and the woman's choice is restricted: in India we have endogamy or, at most, \textbf{restricted hypergamy}; \textbf{hypogamy is prohibited} (marry within your caste or a higher caste, never a lower one).
\item \textbf{Endogamy/hypergamy + control over female sexuality reproduce: (1) caste inequality, (2) the ideology of purity and pollution, and (3) gender subordination.}
\end{itemize}
\end{KeyConcept}

\section{Brahmanical Patriarchy}

\begin{KeyConcept}
\begin{itemize}
\item \textbf{Brahmanical Patriarchy} (a term coined by \textbf{Uma Chakravarti}) = \textbf{Brahmanism + Patriarchy}: Brahmanism legitimises the superiority of the upper-caste Brahmin over the lower Shudra/Dalit; patriarchy legitimises the superiority of men over women. Combined, it means \textbf{upper-caste men ruling over upper-caste women as well as lower-caste women} --- maintaining both the oppressive caste system and patriarchy.
\item The \textbf{bricks} of this structure: \textbf{caste hierarchy} (restricting women's choice; wealth and resources concentrated with the upper castes), \textbf{female purity / stree dharma} (the idea that a woman should strive to be desirable --- which ensures women's participation in their own subordination, since her sexuality is owned by the caste), and \textbf{ritual caste purity} (an elevated spiritual state rewarded for following social norms that oppress women and lower castes).
\item The \textbf{cement} holding these bricks together: the \textbf{Dharma Shastras, Manusmriti and the Yajnavalkya Smriti}.
\end{itemize}
\end{KeyConcept}

\begin{KeyConcept}
\begin{itemize}
\item \textbf{Arranged marriage is exclusive to India} (found nowhere else) --- it exists to maintain caste inequality, by subduing the woman's choice. At the end of the day, \textbf{women become the gateway to caste}.
\item \textbf{Ambedkar} was always more concerned with \textbf{women's education than men's} (upper-caste Brahmins already possess the cultural capital to get educated). \textbf{Nehru} similarly believed women and men are the \textbf{two wheels of the vehicle called India}.
\end{itemize}
\end{KeyConcept}

\begin{QuickRevision}
\begin{itemize}
\item Rudolph \& Rudolph: vertical (upper-caste elites mobilise masses; status quo), horizontal (same-caste mobilises same-caste; Srinivas's horizontal solidarity), differential (political party mobilises across castes) mobilization.
\item Differential = democratization (ideology/development); vertical = no change; conclusion: caste associations = agents of modernity; casteization of politics $\rightarrow$ politicization of caste.
\item Anil Bhatt (four ways): interest \& awareness (high castes higher); identification (no relation to partisanship); influence (middle/low dominant castes). Pre-independence they intersected; now they don't $\rightarrow$ caste democratized.
\item Caste requires endogamy (Surpanakha/Rama; Dharma Shastra; Yajnavalkya Smriti: virgin, same varna, not sapinda/pravara/gotra); caste = social fact; individual is agent of caste reproduction.
\item Uma Chakravarti: marriage + sexuality = foundation of caste; endogamy/hypergamy + control of female sexuality reproduce caste inequality + purity/pollution + gender subordination.
\item Brahmanical Patriarchy = Brahmanism + patriarchy (upper-caste men over women); bricks = caste hierarchy, female purity/stree dharma, ritual caste purity; cement = Dharma Shastra/Manusmriti.
\item Arranged marriage exclusive to India; women = gateway to caste; Ambedkar \& Nehru on women's education.
\end{itemize}
\end{QuickRevision}

